\clearpage
\renewcommand{\sectioncolor}{csD}
\section{Case Study 4 — \texorpdfstring{$e^{\pi i}+1=0$}{e\^{}(pi i)+1=0}}
\markboth{Case Study 4 · How the ideas fit together}{}
\noindent\chip{csD}{BOOK pp.\ 433--452}\;\chip{slate}{PDF pp.\ 452--471}\;\chip{csD}{DELIVERS: why the equation is true \emph{by virtue of what it means}}

\vspace{6pt}
\subsection{Dissolving the pseudo-question}

\begin{critbox}[title={``What does it mean to multiply $e$ by itself $\pi i$ times?''}]
\begin{center}\large\bfseries\color{crimson}
Nothing. The question is malformed.
\end{center}
\begin{bookquote}{book p.~434 / pdf p.~453}
\textit{$e^{\pi i}$ is not an instance of the self-multiplication function. It is, rather, a special
case of the exponential function $w=e^{z}$, where $z=\pi i$.}
\end{bookquote}
\end{critbox}

\noindent$\pi i$ is the \emph{point} $(0,\pi)$ in an input complex plane; the function $w=e^{z}$
sends it to the point $(-1,0)$ in an output complex plane.

\begin{center}
\begin{tikzpicture}[font=\footnotesize]
%---- z-plane ----
\begin{scope}[shift={(0,0)},scale=0.92]
  \fill[csD!5,rounded corners=4pt] (-2.3,-2.3) rectangle (2.3,2.5);
  \draw[slate!50,-{Stealth[length=4pt]}] (-2.05,0)--(2.05,0) node[right,font=\scriptsize]{$x$};
  \draw[slate!50,-{Stealth[length=4pt]}] (0,-2.05)--(0,2.2) node[above,font=\scriptsize]{$yi$};
  \fill[teal] (1.0,0) circle (2.2pt); \node[font=\scriptsize,below right=-3pt,text=teal] at (1.0,0) {$(1,0)$};
  \fill[crimson] (0,1.62) circle (2.4pt);
  \node[font=\scriptsize,right=3pt,text=crimson] at (0,1.62) {$\pi i=(0,\pi)$};
  \fill[ink] (0,0) circle (1.6pt); \node[font=\scriptsize,below left=-3pt] at (0,0) {$0$};
  \node[font=\small\bfseries,text=csD] at (0,-2.62) {the $z$-plane};
\end{scope}
%---- the arrow ----
\draw[-{Stealth[length=8pt,width=7pt]},line width=1.6pt,draw=csD] (2.5,0)--(5.1,0);
\node[font=\small,text=csD,above=2pt] at (3.8,0.08) {$w=e^{z}$};
\node[font=\scriptsize,text=slate,below=2pt,align=center] at (3.8,-0.1)
   {a map from one\\complex plane\\to another};
%---- w-plane ----
\begin{scope}[shift={(7.6,0)},scale=0.92]
  \fill[csD!5,rounded corners=4pt] (-2.3,-2.3) rectangle (2.3,2.5);
  \draw[slate!50,-{Stealth[length=4pt]}] (-2.05,0)--(2.05,0) node[right,font=\scriptsize]{$u$};
  \draw[slate!50,-{Stealth[length=4pt]}] (0,-2.05)--(0,2.2) node[above,font=\scriptsize]{$vi$};
  \fill[teal] (1.45,0) circle (2.2pt); \node[font=\scriptsize,below right=-3pt,text=teal] at (1.45,0) {$(e,0)$};
  \fill[crimson] (-1.0,0) circle (2.4pt);
  \node[font=\scriptsize,below left=-3pt,text=crimson] at (-1.0,0) {$(-1,0)$};
  \fill[ink] (0.55,0) circle (1.6pt); \node[font=\scriptsize,above right=-3pt] at (0.55,0) {$(1,0)$};
  \node[font=\small\bfseries,text=csD] at (0,-2.62) {the $w$-plane};
\end{scope}
%---- annotations ----
\node[anchor=west,font=\scriptsize,text=slate,text width=5.2cm] at (10.35,1.15)
  {$e^{z}$ sends $0\mapsto1$, \ $1\mapsto e$,\\ and --- the equation --- \ $\pi i\mapsto-1$.};
\node[anchor=west,font=\scriptsize,text=crimson,text width=5.2cm] at (10.35,-0.35)
  {\textbf{So $e^{\pi i}=-1$ is a single value of a function.} It is \emph{not} repeated
   multiplication, and it never was.};
\end{tikzpicture}
\end{center}
\vspace{-4pt}
\begin{center}\begin{minipage}{0.93\textwidth}\footnotesize\color{slate}
\textbf{Figure 12.}\; The book's Figure CS\,4.1 \pg{434--435}{453--454}.
\end{minipage}\end{center}

\subsection{Their own anti-climax --- and why it matters}

To their real credit, Lakoff \& Núñez immediately refuse to let the reader be impressed by this.

\begin{bookquote}{book p.~434 / pdf p.~453}
\textit{This ought to leave you cold.\ldots\ Numerically, all that is happening is that
$3.141592654\ldots$ on the $y$-axis in one plane is being mapped by some function onto $-1$ on the
$x$-axis in another\ldots\ From a purely numerical point of view, this has no interest at all!
If all that is at issue is calculating values of functions and seeing which number maps onto which
other number, the greatest equation of all time is a big bore! Euler was too smart to care about
something that boring.}
\end{bookquote}

\begin{keybox}[title={The move that organises the rest of the case study}]
The significance of $e^{\pi i}=-1$ is \textbf{conceptual, not numerical}. It is not given by the
values computed for $w=e^{z}$. It lies in \emph{what the equation tells us about how the branches of
mathematics are related} --- algebra to geometry, geometry to trigonometry, calculus to trigonometry,
and complex arithmetic to all of them \pg{434}{453}.
\end{keybox}

\subsection{The Trigonometric Complex Plane blend}

Case Study 1's \blend{Polar--Trigonometry blend} $\;\oplus\;$ Case Study 3's \blend{$90^{\circ}$
Rotation-Plane blend}. The entailments are the trigonometric form of complex arithmetic:

\begin{center}
\begin{tcolorbox}[enhanced,width=0.95\textwidth,colback=csD!4,colframe=csD,boxrule=1pt,arc=3pt,
  title={\bfseries Entailments in the blend \normalfont\itshape(book p.~437 / pdf p.~456)},
  fonttitle=\bfseries,coltitle=white,colbacktitle=csD]
\begin{center}
$r\cos\theta=a$ \qquad $r(i\sin\theta)=bi$ \qquad
$\underbrace{r\cis\theta}_{\text{\scriptsize notation}}:\;r(\cos\theta+i\sin\theta)=a+bi$
\end{center}
\vspace{2pt}
\[
(r_1\cis\alpha)\cdot(r_2\cis\beta)=(r_1r_2)\cis(\alpha+\beta)
\quad\Longleftrightarrow\quad
(a+bi)(c+di)=(ac-bd)+(ad+bc)i
\]
\end{tcolorbox}
\end{center}

\begin{warmbox}[title={Do not let the familiarity hide the metaphor}]
\begin{bookquote}{book p.~438 / pdf p.~457}
\textit{The complex plane, with complex addition and multiplication, has no inherent trigonometric
structure. Complex arithmetic does not need trigonometry. Trigonometry is in another mathematical
domain altogether. When we understand complex arithmetic in terms of trigonometry, we are
conceptualizing one mathematical domain in terms of another, which is what conceptual metaphor does.}
\end{bookquote}
\end{warmbox}

\clearpage
\subsection{Two pieces of the function \texorpdfstring{$e^{z}$}{e\^{}z}}

Since $e^{z}=e^{x+yi}=e^{x}\cdot e^{yi}$, the function splits into two halves that can be understood
separately and then recombined.

\noindent\begin{minipage}[t]{0.485\textwidth}
\vspace{0pt}
\textbf{\color{csD}(a)\; The $y=0$ slice: $e^{x}$ on the real axis}\\[4pt]
\begin{tikzpicture}[font=\scriptsize]
% z-plane axis
\draw[slate!55,line width=1pt,-{Stealth[length=4pt]}] (-0.3,1.7)--(8.0,1.7);
\foreach \x/\l in {0.6/{$-3$},1.8/{$-2$},3.0/{$-1$},4.2/{$0$},5.4/{$1$},6.6/{$2$},7.4/{$3$}}{
  \fill[slate] (\x,1.7) circle (1.5pt); \node[above=2pt,font=\tiny] at (\x,1.7) {\l};}
\node[anchor=east,font=\tiny\bfseries,text=slate] at (-0.35,1.7) {$z$};
% w-plane axis (log scale look)
\draw[slate!55,line width=1pt,-{Stealth[length=4pt]}] (-0.3,-0.6)--(8.0,-0.6);
\node[anchor=east,font=\tiny\bfseries,text=slate] at (-0.35,-0.6) {$w$};
\fill[ink] (0.6,-0.6) circle (1.5pt); \node[below=2pt,font=\tiny] at (0.6,-0.6) {$0$};
\foreach \x/\l/\c in {3.55/{$1/e$}/teal,4.2/{$1$}/crimson,5.4/{$e$}/csD,6.55/{$e^{2}$}/csD,7.35/{$e^{3}$}/csD}{
  \fill[\c] (\x,-0.6) circle (1.8pt); \node[below=2pt,font=\tiny,text=\c] at (\x,-0.6) {\l};}
% arrows
\foreach \a/\b in {3.0/3.55, 4.2/4.2, 5.4/5.4, 6.6/6.55, 7.4/7.35}{
  \draw[purple!55,-{Stealth[length=3pt]},line width=0.5pt] (\a,1.55)--(\b,-0.42);}
\draw[teal,-{Stealth[length=3pt]},line width=0.5pt] (1.8,1.55)--(3.15,-0.42);
\draw[teal,-{Stealth[length=3pt]},line width=0.5pt] (0.6,1.55)--(2.9,-0.42);
% brackets
\draw[decorate,decoration={brace,amplitude=3.5pt},teal,line width=0.7pt] (0.62,-1.15)--(4.12,-1.15);
\node[font=\tiny,text=teal,below=3pt,text width=3.6cm,align=center] at (2.35,-1.2)
  {the whole negative half-line squeezed into $(0,1)$};
\draw[decorate,decoration={brace,amplitude=3.5pt},csD,line width=0.7pt] (4.3,-1.15)--(7.9,-1.15);
\node[font=\tiny,text=csD,below=3pt,text width=3.4cm,align=center] at (6.1,-1.2)
  {the positive half stretched out to $\infty$};
\node[font=\tiny,text=purple] at (4.2,0.55) {arithmetic $\to$ geometric progressions};
\end{tikzpicture}
\\[26pt]
{\footnotesize\color{slate}Label each image point with its pre-image and you have literally drawn a
\textbf{slide rule} --- the logarithm of Case Study 2 reappearing as a picture \pg{441}{460}, the
book's Figure CS\,4.5.}
\end{minipage}\hfill
\begin{minipage}[t]{0.475\textwidth}
\vspace{0pt}
\textbf{\color{csD}(b)\; The wrapping function $f(\theta)=\cos\theta+i\sin\theta$}\\[4pt]
\begin{center}
\begin{tikzpicture}[font=\scriptsize,scale=0.95]
% segment
\draw[slate!60,line width=1.1pt] (-2.6,2.5)--(2.6,2.5);
\foreach \x/\l/\c in {-2.6/{$-\pi$}/crimson,-1.3/{$-\pi/2$}/purple,0/{$0$}/teal,
                      1.3/{$\pi/2$}/csD,2.6/{$\pi$}/crimson}{
  \fill[\c] (\x,2.5) circle (2pt); \node[above=2pt,font=\tiny,text=\c] at (\x,2.5) {\l};}
\node[anchor=east,font=\tiny,text=slate] at (-2.75,2.5) {$[-\pi,\pi]$};
% circle
\draw[csD,line width=1.2pt] (0,0) circle (1.3);
\draw[slate!35,-{Stealth[length=3pt]}] (-1.75,0)--(1.75,0);
\draw[slate!35,-{Stealth[length=3pt]}] (0,-1.75)--(0,1.75);
\fill[teal]    (1.3,0)  circle (2.2pt); \node[font=\tiny,text=teal,right=2pt]   at (1.3,0)  {$(1,0)$};
\fill[csD]     (0,1.3)  circle (2.2pt); \node[font=\tiny,text=csD,above right=-1pt] at (0,1.3)  {$i$};
\fill[crimson] (-1.3,0) circle (2.4pt); \node[font=\tiny,text=crimson,left=2pt]  at (-1.3,0) {$-1$};
\fill[purple]  (0,-1.3) circle (2.2pt); \node[font=\tiny,text=purple,below right=-1pt] at (0,-1.3) {$-i$};
% wrap arrows
\draw[teal!60,dotted,-{Stealth[length=3.5pt]},line width=0.7pt] (0,2.32)--(1.24,0.10);
\draw[csD!60,dotted,-{Stealth[length=3.5pt]},line width=0.7pt] (1.3,2.32)--(0.12,1.26);
\draw[crimson!60,dotted,-{Stealth[length=3.5pt]},line width=0.7pt]
   (2.6,2.32) .. controls (2.55,1.35) and (-0.4,1.05) .. (-1.24,0.10);
\draw[crimson!60,dotted,-{Stealth[length=3.5pt]},line width=0.7pt]
   (-2.6,2.32) .. controls (-2.55,1.20) .. (-1.36,0.14);
\draw[purple!60,dotted,-{Stealth[length=3.5pt]},line width=0.7pt]
   (-1.3,2.32) .. controls (-1.75,0.55) and (-0.9,-1.05) .. (-0.14,-1.26);
\end{tikzpicture}
\end{center}
\vspace{-2pt}
{\footnotesize\color{slate}The segment is \textbf{wrapped around the unit circle}. Both $\pi$ and
$-\pi$ land on $(-1,0)$ --- that \emph{is} periodicity. Points a distance $\pi$ apart map to points
related by $\times(-1)$; points $\pi/2$ apart, by $\times i$. \textbf{Case Study 3's rotation
metaphor is now doing visible work inside a trigonometric object} \pg{439--440}{458--459}.}
\end{minipage}

\vspace{10pt}
\begin{center}\begin{minipage}{0.93\textwidth}\footnotesize\color{slate}
\textbf{Figure 13.}\; The two halves of $e^{z}=e^{x}\cdot e^{yi}$.
\end{minipage}\end{center}

\clearpage
\subsection{Why \texorpdfstring{$e^{iy}=\cos y+i\sin y$}{e\^{}(iy) = cos y + i sin y}}

The standard route is via Taylor series, and the book runs it honestly: two conceptually different
rules turn out to have the \emph{same derivatives at $0$}, cycling with period four.

\begin{center}\small
\begin{tabular}{@{}c >{\raggedright\arraybackslash}p{4.4cm} >{\raggedright\arraybackslash}p{3.6cm} c >{\raggedright\arraybackslash}p{3.3cm}@{}}
\toprule
\rowcolor{csD!12}
& \multicolumn{2}{l}{\textbf{\color{csD}$f(y)=\cos y+i\sin y$}} &
  \multicolumn{2}{l}{\textbf{\color{csD}$f(y)=e^{yi}$}}\\
\rowcolor{csD!12}
$n$ & derivative & at $y=0$ & derivative & at $y=0$\\
\midrule
0 & $\cos y+i\sin y$   & $1$  & $e^{yi}$        & $1$\\
1 & $-\sin y+i\cos y$  & $i$  & $i\,e^{yi}$     & $i$\\
2 & $-\cos y-i\sin y$  & $-1$ & $-e^{yi}$       & $-1$\\
3 & $\sin y-i\cos y$   & $-i$ & $-i\,e^{yi}$    & $-i$\\
\rowcolor{gold!22}
4 & $\cos y+i\sin y$ \;\textbf{(back to start)} & $1$ & $e^{yi}$ \;\textbf{(back to start)} & $1$\\
\bottomrule
\end{tabular}
\end{center}
\vspace{-2pt}
\begin{center}\begin{minipage}{0.93\textwidth}\footnotesize\color{slate}
\textbf{Table 5.}\; The two derivative cycles \pg{445--446}{464--465}. Same column of values
$\Rightarrow$ same Taylor series $\Rightarrow$ same set of ordered pairs.
\end{minipage}\end{center}

\begin{keybox}[title={What the ``$=$'' in $e^{iy}=\cos y+i\sin y$ actually means}]
This is the cognitive gloss that generalises far beyond Euler's formula. Under the metaphor
\meta{A Function Is a Set of Ordered Pairs}, the equals sign does \textbf{not} say that the two sides
express the same idea.
\begin{center}
\begin{tikzpicture}[font=\footnotesize,
  r/.style={rounded corners=4pt,draw=#1,line width=1pt,fill=#1!8,inner sep=7pt,
            text width=4.6cm,align=center,minimum height=1.5cm}]
 \node[r=csD]  (L) at (-5.0,0) {\textbf{rule 1}\\$e^{yi}$\\[2pt]{\scriptsize an exponential: maps sums to products}};
 \node[r=csA]  (R) at ( 5.0,0) {\textbf{rule 2}\\$\cos y+i\sin y$\\[2pt]{\scriptsize a trigonometric sum on the unit circle}};
 \node[rounded corners=4pt,fill=gold!20,draw=gold!80!black,line width=1pt,inner sep=7pt,
       text width=4.4cm,align=center,minimum height=1.5cm] (S) at (0,0)
       {\textbf{one referent}\\the same \emph{set of ordered pairs}};
 \draw[-{Stealth[length=5pt]},line width=1.1pt,draw=csD] (L)--(S);
 \draw[-{Stealth[length=5pt]},line width=1.1pt,draw=csA] (R)--(S);
\end{tikzpicture}
\end{center}
\vspace{-3pt}
\begin{bookquote}{book p.~446 / pdf p.~465}
\textit{The ``$=$'' means that the two conceptually different mapping rules on the two sides of the
equal sign have the same referent, the same set of ordered pairs.}
\end{bookquote}
\textbf{Worth internalising:} much of the apparent magic of mathematical identities is the collapse of
\emph{same rule} into \emph{same graph}. The book flags this back to Chapter 16, where ``$=$'' is
shown to relate ideas that are cognitively very different.
\end{keybox}

\noindent Setting $y=\pi$ and using $\cos\pi=-1$ (the Trigonometry metaphor, Case Study 1) and
$\sin\pi=0$:
\[
e^{\pi i}=\cos\pi+i\sin\pi=-1+0i=-1
\qquad\Longrightarrow\qquad
\boxed{\;e^{\pi i}+1=0\;}
\]

\clearpage
\subsection{Euler's blend, and the answer that was worth waiting for}

The Taylor-series story tells you \emph{that} the equation holds. It does not tell you \emph{why it
should}. The book's own verdict on its own exposition is disarming:

\begin{bookquote}{book p.~446 / pdf p.~465}
\textit{At first, all this sounds like mumbo-jumbo.\ldots\ Okay. But again, why does it matter that
$e^{y i}=-1$ when $y=\pi$? And what does it mean?}
\end{bookquote}

\noindent The answer comes in two parts: a blend, and then --- finally --- four shared properties.

\subsubsection{Euler's blend: five domains}

\begin{center}
\begin{tikzpicture}[font=\footnotesize,
  d/.style={rounded corners=4pt,draw=csD,line width=1pt,fill=csD!8,inner sep=6pt,
            text width=2.9cm,align=center,minimum height=1.9cm},
  ar/.style={-{Stealth[length=5pt,width=4pt]},line width=1pt,draw=csD!75}]
\node[d] (D1) at (0,0)     {\textbf{1. Power series}\\[3pt]$\sum a_n z^{n}$\\{\scriptsize via the BMI}};
\node[d] (D2) at (3.55,0)  {\textbf{2. Taylor series}\\[3pt]$\sum\frac{f^{(n)}(0)}{n!}z^{n}$\\{\scriptsize via the BMI}};
\node[d] (D3) at (7.10,1.35){\textbf{3. Exponential}\\[3pt]$f(\theta)=e^{i\theta}$};
\node[d] (D4) at (7.10,-1.35){\textbf{4. Trig.\ complex plane}\\[3pt]$f(\theta)=\cos\theta+i\sin\theta$};
\node[d,draw=gold!80!black,fill=gold!14] (D5) at (11.0,0)
   {\textbf{5. Sets of ordered pairs}\\[3pt]$\{(\theta,f(\theta))\}$};
\node[rounded corners=5pt,fill=deepblue,text=white,inner sep=7pt,font=\small\bfseries] (EQ) at (14.9,0)
   {$e^{\pi i}+1=0$};
\draw[ar] (D1)--(D2);
\draw[ar] (D2)--(D3); \draw[ar] (D2)--(D4);
\draw[ar] (D3)--(D5); \draw[ar] (D4)--(D5);
\draw[-{Stealth[length=6pt,width=5pt]},line width=1.3pt,draw=deepblue] (D5)--(EQ);
\node[font=\scriptsize,text=slate,anchor=north,text width=5cm,align=center] at (1.8,-1.15)
   {2 is a special case of 1, so the whole Power-Series blend comes along};
\end{tikzpicture}
\end{center}
\vspace{-2pt}
\begin{center}\begin{minipage}{0.93\textwidth}\footnotesize\color{slate}
\textbf{Figure 14.}\; Euler's blend \pg{447}{466}. The book names it ``partly in honor of Euler and
partly because we suspect that he must have had, at least intuitively and unconsciously, an implicit
understanding of the relationships between the branches of classical mathematics that the blend
represents.''
\end{minipage}\end{center}

\subsubsection{The real answer: four properties, held in common}

\begin{center}
\begin{tikzpicture}[font=\footnotesize,
  fn/.style={rounded corners=4pt,draw=#1,line width=1.2pt,fill=#1!9,inner sep=7pt,
             text width=4.3cm,align=center,minimum height=1.35cm},
  pr/.style={rounded corners=3pt,draw=teal,line width=0.9pt,fill=teal!7,inner sep=6pt,
             text width=6.4cm,align=left,minimum height=1.15cm},
  a1/.style={-{Stealth[length=4.5pt]},line width=0.85pt,draw=csD!70},
  a2/.style={-{Stealth[length=4.5pt]},line width=0.85pt,draw=csA!80}]
\node[fn=csD] (L) at (-5.9,0.0)  {$e^{iy}$\\{\scriptsize an exponential}};
\node[fn=csA] (R) at ( 5.9,0.0)  {$\cos y+i\sin y$\\{\scriptsize a trigonometric sum}};
\node[pr] (p1) at (0,2.55) {\textbf{1.} maps \textbf{sums onto products}\\{\scriptsize $e^{a+b}=e^ae^b$\; ;\; $\cis(\alpha+\beta)=\cis\alpha\cdot\cis\beta$}};
\node[pr] (p2) at (0,0.85) {\textbf{2.} \textbf{changes in exact proportion to its size}\\{\scriptsize entailed by 1, via the factoring step of Case Study 2}};
\node[pr] (p3) at (0,-0.85){\textbf{3.} is \textbf{periodic}\\{\scriptsize \meta{Recurrence Is Circularity} + the unit circle}};
\node[pr,draw=crimson,fill=crimson!6] (p4) at (0,-2.55)
   {\textbf{4.} is \textbf{self-regulating}: $f''=-f$, \;$f''''=f$\\{\scriptsize acceleration is the negative of size}};
\foreach \n in {p1,p2,p3,p4}{\draw[a1] (L)--(\n); \draw[a2] (R)--(\n);}
\end{tikzpicture}
\end{center}
\vspace{-2pt}
\begin{center}\begin{minipage}{0.93\textwidth}\footnotesize\color{slate}
\textbf{Figure 15.}\; The idea analysis of the identity \pg{448--450}{467--469}. Two functions
characterised by \emph{different combinations of concepts} entail \emph{exactly the same properties}.
\end{minipage}\end{center}

\begin{keybox}[title={The sentence the whole book is built to deliver}]
\begin{bookquote}{book p.~450 / pdf p.~469}
\textit{Though they are characterized by different combinations of concepts, they entail, by virtue
of what they mean, exactly the same properties. From the perspective of mathematical idea analysis,
this is why they are equal.}
\end{bookquote}
\end{keybox}

\subsubsection{Self-regulation, made visible}

Property 4 is the one that is easiest to state and hardest to \emph{see}. $f''=-f$ means the
acceleration is the negative of the size: as the function grows, its deceleration grows, until growth
stops and reverses.

\begin{center}
\begin{tikzpicture}
\begin{axis}[width=12.6cm,height=5.0cm,
  axis lines=middle, xmin=-0.4,xmax=13.2,ymin=-1.55,ymax=1.65,
  xtick={0,3.14159,6.28318,9.42477,12.56637},
  xticklabels={$0$,$\pi$,$2\pi$,$3\pi$,$4\pi$},
  ytick={-1,1}, tick label style={font=\tiny}, axis line style={slate!60},
  legend style={font=\tiny,draw=none,fill=none,at={(0.5,1.30)},anchor=north,legend columns=3},
  clip=false]
 \addplot[csD,line width=1.3pt,domain=-0.2:13.1,samples=280]{cos(deg(x))};
 \addlegendentry{$f=\cos y$ \;(size)}
 \addplot[amber,line width=1.1pt,dashed,domain=-0.2:13.1,samples=280]{-sin(deg(x))};
 \addlegendentry{$f'$ \;(rate)}
 \addplot[crimson,line width=1.1pt,dotted,domain=-0.2:13.1,samples=280]{-cos(deg(x))};
 \addlegendentry{$f''=-f$ \;(acceleration)}
 \draw[slate!35,line width=0.4pt] (axis cs:-0.4,0)--(axis cs:13.2,0);
 \node[font=\tiny,text=crimson,align=center] at (axis cs:1.75,-1.32)
   {size positive $\Rightarrow$\\acceleration negative};
 \node[font=\tiny,text=csD,align=center] at (axis cs:4.35,1.32)
   {size negative $\Rightarrow$\\acceleration positive};
 \draw[crimson,line width=0.6pt,-{Stealth[length=3pt]}] (axis cs:1.55,-1.05)--(axis cs:0.95,-0.15);
 \draw[csD,line width=0.6pt,-{Stealth[length=3pt]}] (axis cs:4.35,1.07)--(axis cs:4.2,0.18);
\end{axis}
\end{tikzpicture}
\end{center}
\vspace{-4pt}
\begin{center}\begin{minipage}{0.93\textwidth}\footnotesize\color{slate}
\textbf{Figure 16.}\; ``The self-regulating character comes from the fact that the acceleration is
the negative of the size'' \pg{448}{467}. This is the harmonic oscillator, and the book's
contribution is to give it a \emph{conceptual} rather than merely formal presence inside $e^{i\theta}$.
\end{minipage}\end{center}

\clearpage
\subsection{The philosophical punchline}

\begin{ideabox}[title={Why $e$, $\pi$, $i$, 1 and 0 --- and not 192{,}563{,}947.9853294867}]
\begin{bookquote}{book p.~450--451 / pdf p.~469--470}
\textit{They are not mere numbers; they are the arithmetizations of concepts.\ldots\ Unlike, say,
192,563,947.9853294867, these numbers have conceptual meanings in a system of common, important
nonmathematical concepts, like change, acceleration, recurrence, and self-regulation.}
\end{bookquote}
\end{ideabox}

\begin{center}
\begin{tikzpicture}[font=\footnotesize,
  sym/.style={circle,fill=#1,text=white,font=\large\bfseries,inner sep=6.5pt},
  con/.style={rounded corners=3pt,draw=#1,line width=0.9pt,fill=#1!8,inner sep=6pt,
              text width=3.1cm,align=center,minimum height=1.25cm}]
\node[sym=csB]      (e)  at (0,0)    {$e$};
\node[sym=csA]      (pi) at (3.5,0)  {$\pi$};
\node[sym=csC]      (i)  at (7.0,0)  {$i$};
\node[sym=slate]    (one)at (10.5,0) {$1$};
\node[sym=deepblue] (zer)at (14.0,0) {$0$};
\node[con=csB]      at (0,-1.85)    {\textbf{change}\\{\scriptsize growth proportional to size}};
\node[con=csA]      at (3.5,-1.85)  {\textbf{recurrence}\\{\scriptsize periodicity, half a period}};
\node[con=csC]      at (7.0,-1.85)  {\textbf{rotation}\\{\scriptsize cyclicity, $i^{4}=1$}};
\node[con=slate]    at (10.5,-1.85) {\textbf{multiplicative identity}\\{\scriptsize wholes \& parts}};
\node[con=deepblue] at (14.0,-1.85) {\textbf{additive identity}\\{\scriptsize the empty collection}};
\foreach \n in {e,pi,i,one,zer}{\draw[-{Stealth[length=4pt]},line width=0.8pt,draw=slate!60] (\n)--++(0,-1.05);}
\node[rounded corners=4pt,fill=deepblue,text=white,inner sep=8pt,text width=15.6cm,align=center,
      font=\small] at (7.0,-3.75)
  {These concepts matter in everyday life. \textbf{That} is why these numbers keep turning up ---
   and why the equation that binds them is worth caring about.};
\end{tikzpicture}
\end{center}
\vspace{-2pt}
\begin{center}\begin{minipage}{0.93\textwidth}\footnotesize\color{slate}
\textbf{Figure 17.}\; The closing argument of Part VI \pg{450--451}{469--470}.
\end{minipage}\end{center}

\begin{keybox}[title={The last words of Part VI}]
\begin{bookquote}{book p.~451 / pdf p.~470}
\textit{It is no accident that our branches of mathematics are linked in the way they are. Those
conceptual connections---via conceptual metaphors and blends---express ideas that matter to us. The
way the branches of mathematics are interrelated is a consequence of what is important to us in our
everyday lives and how we conceptualize those concerns. The various branches of mathematics
mathematicize our concerns. The mechanism for doing so is the all too human conceptual metaphor and
blending.}
\end{bookquote}
\end{keybox}
